\documentclass[12pt]{article}
\usepackage[utf-8]{inputenc}
\usepackage{amsmath}
\usepackage{xcolor}
\usepackage{listings}
\usepackage{graphicx}
\usepackage{booktabs}
\usepackage{geometry}
\geometry{margin=1in}
\usepackage{hyperref}

\title{SSIE-500: Homework 5\\
Measuring Complexity of Languages}
\author{Zeinab}
\date{November 10, 2025}

\begin{document}

\maketitle

\section{Introduction}
This report analyzes the complexity of three human languages (English, Spanish, and Arabic) by calculating their information entropy and mutual information between consecutive symbols.

\section{Methodology}

\subsection{Data Sources}
The following copyright-free texts were used for analysis:
\begin{itemize}
    \item \textbf{English}: \textit{Alice's Adventures in Wonderland} by Lewis Carroll (5,658 characters)
    \item \textbf{Spanish}: \textit{Don Quijote} by Miguel de Cervantes (1,386 characters)
    \item \textbf{Arabic}: \textit{Alf Layla wa-Layla} (Arabian Nights - Sindbad the Sailor) (4,643 characters)
\end{itemize}

All texts were obtained from Project Gutenberg and public domain classical literature sources.

\subsection{Formulas}

\subsubsection{Information Entropy}
Information entropy measures the average unpredictability of symbols in a text:

\begin{equation}
H(X) = -\sum_{x \in \mathcal{X}} p(x) \log_2 p(x)
\end{equation}

where $p(x)$ is the probability of symbol $x$ occurring in the text.

\subsubsection{Mutual Information}
Mutual information quantifies the dependency between consecutive symbols:

\begin{equation}
MI(X;Y) = \sum_{x,y} p(x,y) \log_2 \frac{p(x,y)}{p(x)p(y)}
\end{equation}

where $p(x,y)$ is the joint probability of symbols $x$ and $y$ appearing consecutively.

\subsection{Implementation}
The analysis was implemented in Python. The program:
\begin{enumerate}
    \item Reads text files in UTF-8 encoding
    \item Calculates symbol frequencies
    \item Computes information entropy
    \item Analyzes consecutive symbol pairs for mutual information
\end{enumerate}

\section{Results}

\subsection{Summary Table}

\begin{table}[h]
\centering
\begin{tabular}{lcccc}
\toprule
\textbf{Language} & \textbf{Total Chars} & \textbf{Unique Symbols} & \textbf{Entropy (bits)} & \textbf{MI (bits)} \\
\midrule
English & 5,658 & 58 & 4.393104 & 1.120734 \\
Spanish & 1,386 & 44 & 4.275240 & 1.320222 \\
Arabic  & 4,643 & 45 & 4.505785 & 0.969090 \\
\bottomrule
\end{tabular}
\caption{Complexity metrics for three languages}
\label{tab:results}
\end{table}

\subsection{Entropy Comparison}
The entropy values indicate:
\begin{itemize}
    \item \textbf{Arabic} (4.506 bits): Highest entropy, most diverse symbol distribution
    \item \textbf{English} (4.393 bits): Middle entropy
    \item \textbf{Spanish} (4.275 bits): Lowest entropy, most predictable symbol distribution
\end{itemize}

\subsection{Mutual Information Comparison}
The mutual information values show:
\begin{itemize}
    \item \textbf{Spanish} (1.320 bits): Highest MI, strongest sequential dependencies
    \item \textbf{English} (1.121 bits): Moderate MI
    \item \textbf{Arabic} (0.969 bits): Lowest MI, weakest sequential dependencies
\end{itemize}

\section{Discussion}

\subsection{Observations}

\textbf{1. Similar Entropy Values:} All three languages have remarkably similar entropy values (4.28-4.51 bits), suggesting that natural languages tend to maintain similar levels of symbol complexity regardless of alphabet system.

\textbf{2. Spanish's High Mutual Information:} Spanish shows the highest mutual information (1.32 bits) compared to English and Arabic. This suggests strong sequential dependencies in Spanish text, where knowing one character provides significant information about the next.

\textbf{3. Arabic's High Entropy:} Despite having a moderate number of unique symbols (45), Arabic achieves the highest entropy (4.51 bits), indicating a very diverse and relatively unpredictable symbol distribution. This may be due to:
\begin{itemize}
    \item Arabic script's connected writing system where letters change form based on position
    \item Root-based morphology creating diverse symbol combinations
    \item The specific classical text source (Arabian Nights) with rich vocabulary
\end{itemize}

\textbf{4. English Middle Ground:} English (Alice in Wonderland) shows moderate values for both entropy (4.39 bits) and mutual information (1.12 bits), positioning it between Arabic and Spanish in the complexity spectrum.

\subsection{Linguistic Implications}

The results demonstrate that languages achieve similar information density through different mechanisms:
\begin{itemize}
    \item \textbf{Alphabet diversity} (more symbols) vs. \textbf{Sequential structure} (stronger dependencies)
    \item This suggests an underlying optimization in human language for efficient communication
\end{itemize}

\section{Conclusion}

This analysis revealed that while English, Spanish, and Arabic have similar information entropy (ranging from 4.28 to 4.51 bits), they differ in their mutual information characteristics. Spanish exhibits the strongest sequential dependencies (1.32 bits), followed by English (1.12 bits), while Arabic shows the weakest (0.97 bits). Interestingly, Arabic achieves the highest entropy despite having moderate mutual information, suggesting that different writing systems and linguistic structures can achieve similar overall complexity through different organizational principles—some relying more on symbol diversity and others on sequential structure.

\section{Appendix: Complete Python Code}

The complete Python implementation used for this analysis is provided below:

\lstset{
    language=Python,
    basicstyle=\footnotesize\ttfamily,
    keywordstyle=\color{blue},
    commentstyle=\color{green!60!black},
    stringstyle=\color{red},
    showstringspaces=false,
    breaklines=true,
    numbers=left,
    numberstyle=\tiny\color{gray},
    frame=single,
    captionpos=b
}

\begin{lstlisting}[caption={main.py - Language Complexity Analysis Tool}]
import math
from collections import Counter, defaultdict
import os

def calculate_entropy(text):
    """
    Calculate information entropy of the frequency distribution of symbols.
    H(X) = -sum p(x) * log2(p(x))
    """
    if not text:
        return 0.0
    
    # Count frequency of each symbol
    symbol_counts = Counter(text)
    total_symbols = len(text)
    
    # Calculate entropy
    entropy = 0.0
    for count in symbol_counts.values():
        probability = count / total_symbols
        entropy -= probability * math.log2(probability)
    
    return entropy

def calculate_mutual_information(text):
    """
    Calculate mutual information between two consecutive symbols.
    MI(X;Y) = sum p(x,y) * log2(p(x,y) / (p(x) * p(y)))
    """
    if len(text) < 2:
        return 0.0
    
    # Count single symbols
    single_counts = Counter(text)
    total_symbols = len(text)
    
    # Count consecutive symbol pairs
    pair_counts = defaultdict(int)
    for i in range(len(text) - 1):
        pair = (text[i], text[i + 1])
        pair_counts[pair] += 1
    
    total_pairs = len(text) - 1
    
    # Calculate mutual information
    mutual_info = 0.0
    for pair, pair_count in pair_counts.items():
        p_xy = pair_count / total_pairs
        p_x = single_counts[pair[0]] / total_symbols
        p_y = single_counts[pair[1]] / total_symbols
        
        mutual_info += p_xy * math.log2(p_xy / (p_x * p_y))
    
    return mutual_info

def analyze_language(file_path, language_name):
    """
    Analyze a text file and calculate entropy and mutual information.
    """
    try:
        # Try different encodings
        encodings = ['utf-8', 'latin-1', 'cp1252', 'iso-8859-1']
        text = None
        
        for encoding in encodings:
            try:
                with open(file_path, 'r', encoding=encoding) as f:
                    text = f.read()
                break
            except UnicodeDecodeError:
                continue
        
        if text is None:
            print(f"Error: Could not read {file_path} with any encoding")
            return None
        
        # Calculate metrics
        entropy = calculate_entropy(text)
        mutual_info = calculate_mutual_information(text)
        
        # Count statistics
        total_chars = len(text)
        unique_symbols = len(set(text))
        
        result = {
            'language': language_name,
            'file': file_path,
            'total_characters': total_chars,
            'unique_symbols': unique_symbols,
            'entropy': entropy,
            'mutual_information': mutual_info
        }
        
        return result
        
    except FileNotFoundError:
        print(f"Error: File '{file_path}' not found")
        return None
    except Exception as e:
        print(f"Error processing {file_path}: {str(e)}")
        return None

def print_results(results):
    """
    Print analysis results in a formatted way.
    """
    print("\n" + "="*80)
    print("LANGUAGE COMPLEXITY ANALYSIS RESULTS")
    print("="*80 + "\n")
    
    for result in results:
        if result:
            print(f"Language: {result['language']}")
            print(f"File: {result['file']}")
            print(f"Total Characters: {result['total_characters']:,}")
            print(f"Unique Symbols: {result['unique_symbols']}")
            print(f"Information Entropy: {result['entropy']:.6f} bits")
            print(f"Mutual Information: {result['mutual_information']:.6f} bits")
            print("-" * 80 + "\n")
    
    # Comparison section
    if len(results) >= 2:
        print("\n" + "="*80)
        print("COMPARISON")
        print("="*80 + "\n")
        
        # Sort by entropy
        sorted_by_entropy = sorted([r for r in results if r], 
                                   key=lambda x: x['entropy'], reverse=True)
        print("Languages ranked by Entropy (highest to lowest):")
        for i, result in enumerate(sorted_by_entropy, 1):
            print(f"{i}. {result['language']}: {result['entropy']:.6f} bits")
        
        print("\nLanguages ranked by Mutual Information (highest to lowest):")
        sorted_by_mi = sorted([r for r in results if r], 
                             key=lambda x: x['mutual_information'], reverse=True)
        for i, result in enumerate(sorted_by_mi, 1):
            print(f"{i}. {result['language']}: {result['mutual_information']:.6f} bits")
        
        print("\n" + "="*80)

def main():
    """
    Main function to analyze text files for different languages.
    """
    print("Language Complexity Analysis")
    print("SSIE-500: Homework 5\n")
    
    # Define the text files for each language
    languages = [
        ('english.txt', 'English'),
        ('spanish.txt', 'Spanish'),
        ('arabic.txt', 'Arabic')
    ]
    
    results = []
    
    print("Analyzing text files...\n")
    for file_path, language_name in languages:
        if os.path.exists(file_path):
            print(f"Processing {language_name}...")
            result = analyze_language(file_path, language_name)
            results.append(result)
        else:
            print(f"Warning: {file_path} not found.")
            results.append(None)
    
    print_results(results)

if __name__ == "__main__":
    main()
\end{lstlisting}

\end{document}
